
\documentclass{article}
\usepackage{colortbl}
\usepackage{makecell}
\usepackage{multirow}
\usepackage{supertabular}

\begin{document}

\newcounter{utterance}

\centering \large Interaction Transcript for game `cladder', experiment `full\_v1.5\_default', episode 2287 with reflexion{-}Qwen3.5{-}4B.
\vspace{24pt}

{ \footnotesize  \setcounter{utterance}{1}
\setlength{\tabcolsep}{0pt}
\begin{supertabular}{c@{$\;$}|p{.15\linewidth}@{}p{.15\linewidth}p{.15\linewidth}p{.15\linewidth}p{.15\linewidth}p{.15\linewidth}}
    \# & \multicolumn{2}{c}{Player} && \multicolumn{2}{c}{Game Master} \\
    \hline

    \theutterance \stepcounter{utterance}  
    & & & \multicolumn{4}{p{0.6\linewidth}}{
        \cellcolor[rgb]{0.9,0.9,0.9}{
            \makecell[{{p{\linewidth}}}]{
                \texttt{\tiny{[P1$\langle$GM]}}
                \texttt{You have attempted this same task before and did not succeed. These are the plans you wrote after those attempts.} \\
\texttt{Trial \#1: New plan: I failed because I calculated the conditional probabilities correctly but ignored the strict formatting constraint to start the answer with "Yes" or "No" before providing the reasoning. Although the logical answer was "No", I initially wrote "Yes" and then tried to correct myself mid{-}stream, which likely violated the "start with" rule or the expectation of a definitive initial stance. Next time, I will first perform the calculation to determine the correct boolean answer (Yes/No), and then immediately output that single word as the very first characters of my response, followed by the reasoning, ensuring no contradictory opening is made.} \\
\\ 
\\ 
\texttt{You are an expert in causal inference. The following question is not a typical commonsense query, but rather a meticulously designed question created by a professor specializing in causal inference, intended to assess the students' mastery of the course content.} \\
\texttt{Imagine a self{-}contained, hypothetical world with only the following conditions, and without any unmentioned factors or causal relationships: Smoking has a direct effect on effort and college admission. Effort has a direct effect on college admission. The overall probability of smoking is 93\%. The probability of nonsmoking and college admission is 3\%. The probability of smoking and college admission is 62\%. Is the chance of college admission smaller when observing smoking?} \\
\texttt{Start your answer with 'Yes' or 'No', followed by additional reasoning or evidence to support your explanation.} \\
            }
        }
    }
     \\ \\

    \theutterance \stepcounter{utterance}  
    & \multicolumn{4}{p{0.6\linewidth}}{
        \cellcolor[rgb]{0.561,0.737,0.561}{
            \makecell[{{p{\linewidth}}}]{
                \texttt{\tiny{[P1$\rangle$GM]}}
                \texttt{No. To determine whether the chance of college admission is smaller when observing smoking, we must compare the conditional probability of admission given smoking (\$P(\textbackslash{}text\{Admission\}|\textbackslash{}text\{Smoking\})\$) with the overall probability of admission (\$P(\textbackslash{}text\{Admission\})\$).} \\
\\ 
\texttt{First, let's define the events:} \\
\texttt{*   \$S\$: Smoking} \\
\texttt{*   \$\textbackslash{}neg S\$: Nonsmoking} \\
\texttt{*   \$A\$: College Admission} \\
\texttt{*   \$\textbackslash{}neg A\$: No College Admission} \\
\\ 
\texttt{We are given the following marginal probabilities:} \\
\texttt{*   \$P(S) = 0.93\$} \\
\texttt{*   \$P(\textbackslash{}neg S) = 1 {-} 0.93 = 0.07\$} \\
\texttt{*   \$P(A \textbackslash{}cap \textbackslash{}neg S) = 0.03\$ (Probability of nonsmoking and admission)} \\
\texttt{*   \$P(A \textbackslash{}cap S) = 0.62\$ (Probability of smoking and admission)} \\
\\ 
\texttt{Next, we calculate the overall probability of college admission, \$P(A)\$, using the law of total probability:} \\
\texttt{\$\$P(A) = P(A \textbackslash{}cap S) + P(A \textbackslash{}cap \textbackslash{}neg S)\$\$} \\
\texttt{\$\$P(A) = 0.62 + 0.03 = 0.65\$\$} \\
\\ 
\texttt{Now, we calculate the conditional probability of admission given smoking, \$P(A|S)\$:} \\
\texttt{\$\$P(A|S) = \textbackslash{}frac\{P(A \textbackslash{}cap S)\}\{P(S)\}\$\$} \\
\texttt{\$\$P(A|S) = \textbackslash{}frac\{0.62\}\{0.93\} \textbackslash{}approx 0.6667\$\$} \\
\\ 
\texttt{Finally, we compare the two probabilities:} \\
\texttt{*   Overall probability of admission: \$0.65\$} \\
\texttt{*   Probability of admission given smoking: \$\textbackslash{}approx 0.6667\$} \\
\\ 
\texttt{Since \$0.6667 > 0.65\$, the chance of college admission is actually **larger** when observing smoking, not smaller. Therefore, the statement "the chance of college admission is smaller when observing smoking" is false.} \\
            }
        }
    }
    & & \\ \\

    \theutterance \stepcounter{utterance}  
    & & & \multicolumn{2}{p{0.3\linewidth}}{
        \cellcolor[rgb]{0.9,0.9,0.9}{
            \makecell[{{p{\linewidth}}}]{
                \texttt{\tiny{[GM$|$GM]}}
                \texttt{No} \\
            }
        }
    }
    & & \\ \\

    \theutterance \stepcounter{utterance}  
    & & & \multicolumn{2}{p{0.3\linewidth}}{
        \cellcolor[rgb]{0.9,0.9,0.9}{
            \makecell[{{p{\linewidth}}}]{
                \texttt{\tiny{[GM$|$GM]}}
                \texttt{no} \\
            }
        }
    }
    & & \\ \\

    \theutterance \stepcounter{utterance}  
    & & & \multicolumn{2}{p{0.3\linewidth}}{
        \cellcolor[rgb]{0.9,0.9,0.9}{
            \makecell[{{p{\linewidth}}}]{
                \texttt{\tiny{[GM$|$GM]}}
                \texttt{game\_result = WIN} \\
            }
        }
    }
    & & \\ \\

\end{supertabular}
}

\end{document}
